\section{Bulgular}
\label{sec:bulgular}

Bu bölüm, sınır sürekliliği ve hesaplama maliyeti bulgularını, ardından
dört ablasyon analizinin sonuçlarını ve son olarak gerçek Isaac Sim
eğitim karşılaştırmasının sonuçlarını sunmaktadır. Yöntem tekrarı
yapılmamış, veri/deney tasarımı ayrıntıları \cref{sec:yontem}'e
bırakılmıştır.

\subsection{Sınır Sürekliliği ve Hesaplama Maliyeti}
\label{sec:sureklilik}

\Cref{tab:continuity}, otuz tohum üzerinde toplanmış (aggregated),
yalnızca söndürülmüş ince detay bileşenini izole eden başlık (headline)
süreklilik ölçümlerini, smoothstep'e göre oran olarak özetlemektedir.
Üretim öncesi Gauss (\texttt{gaussian\_shipped}, $k=1{,}0$), smoothstep'e
kıyasla ortalama 3192 kat daha kötü C0 (yükseklik) sürekliliği ve 535 kat
daha kötü C1 (eğim) sürekliliği üretmektedir; her iki fark da Wilcoxon
işaretli sıra testiyle son derece anlamlıdır ($p=1{,}86\times10^{-9}$,
$n=30$). Kalibre edilmiş Gauss ($k=\ln(1000)$), bu kusuru büyük ölçüde
gidermekte, ancak yine de smoothstep'e göre ölçülebilir düzeyde (C0'da
yaklaşık 9,3 kat, C1'de yaklaşık 1,6 kat) daha yüksek kalıntı hata
bırakmaktadır. Hibrit yöntem, $r=R$ sınırında analitik olarak smoothstep
ile birebir aynı olduğundan ($\mathrm{switch}(R)=0$), otuz tohumun
tamamında sıfır fark göstermiştir; bu, bir ölçüm sonucu değil, yöntemin
tanımından kaynaklanan beklenen bir sonuçtur.

\begin{table}[t]
\centering
\caption{Otuz tohum üzerinde toplanmış, yalnızca ince-detay bileşenini
izole eden sınır sürekliliği oranları (smoothstep'e göre).}
\label{tab:continuity}
\footnotesize
\begin{tabular}{@{}lccc c@{}}
\toprule
\textbf{Method} & \textbf{C0 oranı} & \textbf{C1 oranı} & \textbf{Wilcoxon $p$} & \textbf{$n$ (seed)} \\
\midrule
smoothstep (referans) & 1.0 & 1.0 & -- & 30 \\
gaussian\_shipped ($k$=1.0) & 3192x & 535x & 1.86e-09 & 30 \\
gaussian\_calibrated ($k$=ln(1000)) & 9.3x & 1.6x & 1.86e-09 & 30 \\
hybrid & 1.0 (özdeş) & 1.0 (özdeş) & n/a (analitik olarak özdeş) & 30 \\
\bottomrule
\end{tabular}
\end{table}

\Cref{fig:radial-profile}, dört yöntem varyantının birleşik (makro + ince
detay) yükseklik kesitini göstermektedir. Bu ölçekte dört eğri neredeyse
tam olarak üst üste binmektedir, çünkü ince detay genliği (santimetre
mertebesinde) makro arazi ölçeğine (metre mertebesinde) kıyasla çok
küçüktür; \cref{tab:continuity}'de raporlanan büyüklük mertebesindeki
farklar bu görsel ölçekte çıplak gözle seçilemez. Bu,
\cref{sec:analitikprotokol}'te tanımlanan ``detail-only'' izolasyon
adımının neden gerekli olduğunu doğrudan göstermektedir: dikiş kusuru,
yalnızca ince detay bileşeni ayrıştırıldığında ortaya çıkmaktadır.

\begin{figure}[t]
\centering
\includegraphics[width=\linewidth]{fig04_radial_profile.png}
\caption{Üç karıştırma yönteminin radyal söndürme profili.}
\label{fig:radial-profile}
\end{figure}

\Cref{tab:cost}, yalnızca söndürme formülünün kendisinin (arazi
örneklemesinden izole) mikro-saniye cinsinden hesaplama maliyetini
göstermektedir. Smoothstep en düşük maliyete sahiptir (93,1
mikrosaniye); Gauss varyantları (üretim öncesi ve kalibre edilmiş)
benzer, orta düzey bir maliyete sahiptir (yaklaşık 149--150 mikrosaniye);
hibrit yöntem ise en yüksek maliyete sahiptir (407,8 mikrosaniye,
smoothstep'in yaklaşık 4,3 katı), çünkü her sorguda hem Gauss hem
smoothstep dallarını hem de switch ağırlığını hesaplamaktadır.

\begin{table}[t]
\centering
\caption{Söndürme (fade) formülünün sorgu başına hesaplama maliyeti.}
\label{tab:cost}
\begin{tabular}{@{}lc@{}}
\toprule
\textbf{Method} & \textbf{Cost (\si{\micro\second}/query)} \\
\midrule
smoothstep & 93.12 \\
gaussian\_shipped & 148.84 \\
gaussian\_calibrated & 149.61 \\
hybrid & 407.83 \\
\bottomrule
\end{tabular}
\end{table}

\Cref{fig:continuity-bars}, \cref{tab:continuity}'deki analitik ve
mesh-track süreklilik ölçümlerini yöntem başına çubuk grafik olarak
görselleştirmektedir; \texttt{gaussian\_shipped}'in dört panelin
tamamında diğer üç yönteme kıyasla düzenlerce daha yüksek hata çubuğuna
sahip olması, tabloda raporlanan büyüklük farkını görsel olarak da
doğrulamaktadır.

\begin{figure}[t]
\centering
\includegraphics[width=\linewidth]{fig05_continuity_bars.png}
\caption{Yöntemler arası C0/C1 sınır sürekliliği karşılaştırması.}
\label{fig:continuity-bars}
\end{figure}

\Cref{tab:vram}, iki mesh muhasebe biçimi için bellek tahminini
göstermektedir; bu tahmin yalnızca ağ geometrisine (vertex sayısı)
bağlıdır ve karıştırma yönteminin şekli VRAM üzerinde ölçülebilir bir
etki yaratmamaktadır -gerçek VRAM kaldıracı, karıştırma fonksiyonu
seçimi değil, ağ topolojisi tercihidir (iki-mesh'e karşı tek-grid).

\begin{table}[t]
\centering
\caption{Ağ muhasebe biçimine göre tahmini VRAM kullanımı.}
\label{tab:vram}
\begin{tabular}{@{}lc@{}}
\toprule
\textbf{Mesh accounting scheme} & \textbf{Estimated memory (MB)} \\
\midrule
Two-mesh (literal: separate global + local meshes) & 5.24 \\
Single-grid (baked: one combined mesh) & 1.05 \\
\bottomrule
\end{tabular}
\end{table}

\Cref{fig:cost-comparison}, \cref{tab:cost} ve \cref{tab:vram}'teki maliyet
ölçütlerini tek bir görünümde birleştirmektedir: mesh-track üretim
süresi ve analytic-track sorgu verimi yöntemler arasında görece
yakınken, hibrit'in fade-only maliyetteki belirgin sıçraması ve dört
yöntemin de aynı VRAM ayak izini paylaştığı aynı grafik takımında bir
arada görülebilmektedir.

\begin{figure}[t]
\centering
\includegraphics[width=\linewidth]{fig06_cost_comparison.png}
\caption{Yöntemler arası hesaplama maliyeti karşılaştırması.}
\label{fig:cost-comparison}
\end{figure}

\subsection{Ablasyon Analizleri}
\label{sec:ablasyon}

Bu bölüm, önerilen hibrit yöntemin dört tasarım parametresine (temas
yarıçapı, switch genişliği, geçiş yarıçapı, grid çözünürlüğü) karşı
sağlamlığını sınayan dört ayrı ablasyon analizini sunmaktadır.

Hibrit yöntemin temas yarıçapı ($r_c$) taramasında (1,0--4,0~m;
\cref{fig:ablation-rc}), switch'in kendi iç geçiş sınırındaki ($r=r_c$)
C0 süreksizliği tüm $r_c$ değerlerinde $3\times10^{-4}$~m'nin altında
kalmış ve söndürme maliyeti $r_c$'den bağımsız olarak sabit (yaklaşık
755--793 mikrosaniye) kalmıştır; bu, maliyetin switch'in nerede
konumlandığından değil, iki dalın da her durumda hesaplanmasından
kaynaklandığını doğrulamaktadır.

\begin{figure}[t]
\centering
\includegraphics[width=\linewidth]{fig07_ablation_rc.png}
\caption{Temas yarıçapı ($r_c$) taraması: switch iç sürekliliği ve
maliyet.}
\label{fig:ablation-rc}
\end{figure}

Switch genişliği ($w$) taramasında (0,05--2,0~m; \cref{fig:ablation-w}),
switch'in iç sürekliliği $w=0{,}05$~m'de $7{,}03\times10^{-4}$~m'den
$w=1{,}0$~m'de $1{,}74\times10^{-7}$~m'ye kadar düzenli olarak
iyileşmiş, ancak $w=2{,}0$~m'de hafifçe kötüleşmiştir
($1{,}83\times10^{-5}$~m); bu, switch genişliğinin temas yarıçapına yakın
büyüklüklere ulaştığında switch'in yerel eğriliğini bozmaya başladığını
düşündürmekte ve makine hassasiyetine yakın bir iç süreklilik için $w$
yaklaşık 0,3--1,0~m aralığının uygun bir varsayılan olduğunu
göstermektedir; üretim yapılandırmasında kullanılan $w=0{,}3$~m, bu
aralığın içinde bilinçli bir güvenlik marjı sağlamaktadır.

\begin{figure}[t]
\centering
\includegraphics[width=\linewidth]{fig08_ablation_w.png}
\caption{Switch genişliği ($w$) taraması: iç süreklilik.}
\label{fig:ablation-w}
\end{figure}

Bu kusurun geçiş yarıçapından bağımsız olup olmadığını sınamak amacıyla
yapılan yarıçap taramasında ($R=4,6,8,12,16$~m; \cref{fig:ablation-R}),
üretim öncesi Gauss'un C0 (RMS) hatası neredeyse sabit kalmaktadır
(yaklaşık $1{,}68\times10^{-3}$ ile $1{,}76\times10^{-3}$~m arası), yani
hata geçiş yarıçapı büyüdükçe küçülmemektedir; buna karşın smoothstep ve
kalibre edilmiş Gauss'un hatası, yarıçap büyüdükçe (eğrilik azaldıkça)
beklendiği gibi düzenli olarak azalmaktadır. Bu, üretim öncesi Gauss'un
kusurunun bir ayrıklaştırma (discretization) hatası değil, geçiş
yarıçapından bağımsız, ölçek değişmez bir yapısal özellik olduğunu
göstermektedir.

\begin{figure}[t]
\centering
\includegraphics[width=\linewidth]{fig09_ablation_R.png}
\caption{Geçiş yarıçapı ($R$) taraması: ölçek değişmezliği.}
\label{fig:ablation-R}
\end{figure}

Grid çözünürlüğü taramasında (grid aralığı yaklaşık 0,31--1,27~m;
\cref{fig:ablation-grid}), smoothstep ve hibrit yöntemlerin ampirik
yakınsama derecesi (log-log doğrusal regresyon eğimi) yaklaşık 0,90
olarak ölçülmüş, yani hata çözünürlük arttıkça beklenen şekilde
azalmaktadır. Kalibre edilmiş Gauss'ta ise negatif bir eğim (yaklaşık
$-0{,}36$) gözlenmiştir; ancak bu bir fiziksel ``çözünürlük arttıkça
kötüleşme'' bulgusu değildir -kalibre edilmiş Gauss'un hatası zaten
$10^{-6}$ ile $10^{-7}$~m mertebesinde, kalibrasyon sabitinin
belirlediği sabit bir kesme-kalıntısı gürültü tabanına yakındır ve bu
taban, grid çözünürlüğünden bağımsız bir ölçüm sınırı oluşturmaktadır; bu
husus, yanıltıcı bir fiziksel yorumdan kaçınmak için burada açıkça
belirtilmektedir.

\begin{figure}[t]
\centering
\includegraphics[width=\linewidth]{fig10_ablation_grid.png}
\caption{Grid çözünürlüğü taraması: ampirik yakınsama derecesi.}
\label{fig:ablation-grid}
\end{figure}

Dört analizin ortak sonucu, önerilen hibrit yöntemin süreklilik
avantajının, makul bir parametre aralığında ($r_c=1{,}0$--$4{,}0$~m,
$w=0{,}3$--$1{,}0$~m) sağlam (robust) olduğu, yani belirli bir ince
ayara aşırı duyarlı olmadığıdır. Tek istisna, çok dar switch
genişliklerinde ($w\le 0{,}1$~m) iç sürekliliğin gözle görülür biçimde
bozulmasıdır; bu, üretim yapılandırmasında ($w=0{,}3$~m) kullanılan
varsayılanın bilinçli bir güvenlik marjı içerdiğini bir kez daha
doğrulamaktadır.

\subsection{Gerçek Isaac Sim Eğitim Sonuçları}
\label{sec:egitimsonuclari}

Dokuz eğitim koşusunun (3 yöntem x 3 tohum) tamamı hatasız
tamamlanmış, her biri 5.734.400 ortam adımına ulaşmıştır
(\cref{tab:training}). Duvar saati (wall-clock) süresi, yöntem
ortalaması olarak smoothstep için $2377\pm 98$~s, kalibre edilmiş Gauss
için $2850\pm 898$~s ve hibrit için $2674\pm 599$~s'dir. Bu ortalamalar,
ilk (seed=42) koşularda gözlenen büyük farkın (Gauss ve hibrit sırasıyla
\%56 ve \%35 daha yavaş) çoklu tohumla büyük ölçüde küçüldüğünü
göstermektedir: seed=7 ve seed=123 koşularının üçü de birbirine yakın
(2318--2345~s) sürerken, yalnızca seed=42'deki Gauss ve hibrit koşuları
belirgin biçimde daha yavaştır. Standart sapmaların büyüklüğü (898 ve
599~s, ortalamanın \%20--30'u) bu tek-koşu sapmasının kaynağıdır;
dolayısıyla ilk ölçümdeki ``\%35--56 daha yavaş'' bulgusu, karıştırma
formülünün kendisinden değil, o belirli koşuya özgü sistem-düzeyi bir
değişkenlikten (arka plan yükü, termal durum veya terrain havuzu
zamanlaması) kaynaklanmış olabilir; üç tohumun ortalamasına bakıldığında
Gauss ve hibrit hâlâ smoothstep'ten yavaş görünmektedir (+\%20 ve
+\%12,5) ama bu fark artık ilk ölçümdeki kadar keskin değildir.

\begin{table}[t]
\centering
\caption{Gerçek Isaac Sim eğitiminde son yüzde onluk pencere özeti
(yöntem başına üç tohum).}
\label{tab:training}
\footnotesize
\begin{tabular}{@{}lccccc@{}}
\toprule
\textbf{Method} & \textbf{$n$} & \textbf{Wall time (s)} & \textbf{Success rate} & \textbf{Timeout rate} & \textbf{Harsh-landing} \\
\midrule
smoothstep & 3 & $2377\pm 98$ & $0.278\pm 0.055$ & $0.499\pm 0.059$ & $0.220\pm 0.022$ \\
gaussian (cal.) & 3 & $2850\pm 898$ & $0.325\pm 0.051$ & $0.454\pm 0.074$ & $0.218\pm 0.038$ \\
hybrid & 3 & $2674\pm 599$ & $0.268\pm 0.025$ & $0.515\pm 0.019$ & $0.214\pm 0.007$ \\
\bottomrule
\end{tabular}
\end{table}

Son yüzde on penceresinde ölçülen nihai başarı oranı, yöntem ortalaması
olarak smoothstep için \%27,8~$\pm$~\%5,5, kalibre edilmiş Gauss için
\%32,5~$\pm$~\%5,1 ve hibrit için \%26,8~$\pm$~\%2,5 olarak bulunmuştur
(\cref{tab:training}). Gauss'un ortalaması hâlâ en yüksektir, ancak üç
yöntemin standart sapmaları birbirleriyle örtüşmektedir. \Cref{tab:pairwise},
bu farkların ikili istatistiksel testlerini özetlemektedir: hiçbir ikili
karşılaştırma istatistiksel olarak anlamlı değildir; başarı oranında
smoothstep-Gauss $p=0{,}20$, smoothstep-hibrit $p=1{,}00$, Gauss-hibrit
$p=0{,}40$; zaman aşımı oranında sırasıyla $p=0{,}70$, $p=1{,}00$,
$p=0{,}40$; sert iniş oranında üçü de $p\ge 0{,}70$. Bootstrap \%95 güven
aralıklarının tamamı sıfırı kapsamaktadır (ör.\ Gauss-hibrit başarı
oranı farkı için $[-0{,}000,\,0{,}104]$ -sınırda ama sıfırı içeriyor).

\begin{table}[t]
\centering
\caption{Son yüzde onluk pencere üzerinde ikili karşılaştırmalar
(Mann-Whitney U, betimsel; bkz. \cref{sec:egitimprotokolu} çekincesi).}
\label{tab:pairwise}
\footnotesize
\begin{tabular}{@{}llcccc@{}}
\toprule
\textbf{Metric} & \textbf{Comparison} & \textbf{$n$} & \textbf{Mean diff} & \textbf{95\% CI} & \textbf{MW $p$} \\
\midrule
Success rate & smoothstep vs gaussian & 3 vs 3 & $-0.047$ & $[-0.112, 0.018]$ & 2.00e-01 \\
Success rate & smoothstep vs hybrid & 3 vs 3 & $0.010$ & $[-0.040, 0.068]$ & 1.00e+00 \\
Success rate & gaussian vs hybrid & 3 vs 3 & $0.057$ & $[-0.000, 0.104]$ & 4.00e-01 \\
Timeout rate & smoothstep vs gaussian & 3 vs 3 & $0.045$ & $[-0.044, 0.133]$ & 7.00e-01 \\
Timeout rate & smoothstep vs hybrid & 3 vs 3 & $-0.017$ & $[-0.073, 0.039]$ & 1.00e+00 \\
Timeout rate & gaussian vs hybrid & 3 vs 3 & $-0.061$ & $[-0.132, 0.008]$ & 4.00e-01 \\
Harsh-landing & smoothstep vs gaussian & 3 vs 3 & $0.002$ & $[-0.041, 0.040]$ & 1.00e+00 \\
Harsh-landing & smoothstep vs hybrid & 3 vs 3 & $0.006$ & $[-0.012, 0.031]$ & 1.00e+00 \\
Harsh-landing & gaussian vs hybrid & 3 vs 3 & $0.004$ & $[-0.027, 0.047]$ & 7.00e-01 \\
\bottomrule
\end{tabular}
\end{table}

\Cref{fig:training-curves}, üç yöntemin dokuz koşusunun tamamını üç ana
sonuç metriği (başarı, zaman aşımı, sert iniş oranı) için aynı grafik
üzerinde göstermektedir; dokuz eğrinin de eğitim boyunca büyük ölçüde iç
içe geçmiş olması, yöntemler arası farkın tohumlar arası varyanstan
görsel olarak ayırt edilemediğini doğrulamaktadır. \Cref{fig:training-diag},
aynı dokuz koşu için müfredat zorluğu, entropi katsayısı ve
xy-ilerleme ödülü gibi tamamlayıcı eğitim diyagnostiklerini
göstermektedir; entropi katsayısının üç yöntemde de neredeyse özdeş bir
yakınsama eğrisi izlemesi, SAC'ın keşif-sömürü dengesinin karıştırma
yöntemi seçiminden bağımsız işlediğine işaret etmektedir.
\Cref{fig:final-window-box}, son yüzde onluk pencere değerlerini yöntem
başına dağılım noktaları ve ortalama $\pm$ standart sapma çubuğu olarak
sunmaktadır; üç yöntemin hata çubuklarının geniş ölçüde örtüştüğü buradan
da açıkça görülmektedir.

\begin{figure}[t]
\centering
\includegraphics[width=\linewidth]{fig11_training_curves.png}
\caption{Üç yöntem için eğitim boyunca izlenen sonuç metrikleri
(başarı/zaman aşımı/sert iniş oranı).}
\label{fig:training-curves}
\end{figure}

\begin{figure}[t]
\centering
\includegraphics[width=\linewidth]{fig12_training_diagnostics.png}
\caption{Üç yöntem için eğitim boyunca izlenen tamamlayıcı
diyagnostikler (müfredat zorluğu, entropi katsayısı, xy-ilerleme ödülü).}
\label{fig:training-diag}
\end{figure}

\begin{figure}[t]
\centering
\includegraphics[width=\linewidth]{fig13_final_window_box.png}
\caption{Son yüzde onluk pencere dağılımları (kutu grafiği).}
\label{fig:final-window-box}
\end{figure}

Bu sonuç, ilk (tek tohumlu) ölçümde gözlenen ve \cref{sec:tartisma}'te
ayrıntılı olarak tartışılan ``Gauss ve hibrit'in smoothstep'ten anlamlı
derecede daha başarılı olduğu'' bulgusunu doğrulamamaktadır. $n=1$'den
$n=3$'e çıkıldığında etki büyüklüğü küçülmüş (Gauss-smoothstep farkı
\%11'den \%4,7'ye düşmüş) ve yön kısmen değişmiştir
(hibrit-smoothstep farkı işaret değiştirmiştir: ilk ölçümde hibrit
smoothstep'ten yüksekti, üç-tohum ortalamasında hibrit smoothstep'e
neredeyse eşittir, hatta hafifçe düşüktür). Bu örüntü, ilk bulgunun -en
azından büyük ölçüde- rastgele koşu-koşu varyansıyla tutarlı olduğunu
göstermektedir; bu husus \cref{sec:tartisma}'te ayrıntılı olarak
tartışılmaktadır.

Sample-efficiency ölçütlerinde de benzer bir görünüm vardır: AUC (eğitim
boyunca ortalama başarı oranı) yöntem başına üç tohum arasında geniş bir
aralıkta değişmektedir (smoothstep 0,381--0,395, Gauss 0,366--0,402,
hibrid 0,339--0,399) ve yöntemler arası fark, tohumlar arası farktan
belirgin biçimde küçüktür -bu da yine tohum varyansının baskın kaynak
olduğuna işaret etmektedir.

Müfredat zorluğu, üç yöntemde de ortalama olarak birbirine yakındır
(smoothstep $0{,}276\pm 0{,}025$, Gauss $0{,}250\pm 0{,}000$, hibrit
$0{,}260\pm 0{,}017$); Gauss'un üç tohumda da tam olarak $0{,}25$'te
sabit kalması (sıfır standart sapma) dikkat çekicidir ve rolling-window
başarı eşiğinin bu üç koşuda da aynı basamakta durduğunu göstermektedir
-bu, ilk ölçümde öne sürülen ``Gauss'un yüksek ortalama başarıya rağmen
müfredatta geride kaldığı'' gözlemiyle tutarlıdır, ancak nedeni doğrudan
test edilmemiştir.
